\documentclass[a4paper,12pt]{article} % Imposta il tipo di documento, formato carta e dimensione base del font

% --- GESTIONE LINGUE E CARATTERI ---
\usepackage[LGR,T1]{fontenc}      % Mappa i caratteri in uscita (T1 per l'italiano/europeo, LGR per il greco)
\usepackage[utf8]{inputenc}       % Gestisce la codifica del file sorgente (permette di scrivere direttamente à, è, ì)
\usepackage[greek,italian]{babel} % Adatta sillabazione e nomi predefiniti (es. "Capitolo") alle lingue scelte
\usepackage{textalpha}            % Permette di inserire lettere greche direttamente nel testo, senza ambiente matematico

% --- MATEMATICA E FONT ---
\usepackage{amsmath}              % Fornisce strumenti avanzati per scrivere formule matematiche
\usepackage{newtxtext,newtxmath}  % 1.  Cambia il font del documento e della matematica in stile "Times"

% 2.    DOPO diciamo a LaTeX di usare il Computer Modern (cmr) per il greco
%       al posto della famiglia di newtxtext (ntxtlf) che non ha quei glifi
\DeclareFontFamilySubstitution{LGR}{ntxtlf}{cmr}

\usepackage{xcolor}               % Gestione dei colori
\definecolor{darkred}{rgb}{0.6, 0, 0} % rosso scuro

% --- LAYOUT E FORMATTAZIONE ---
\usepackage{setspace}             % Permette di modificare facilmente l'interlinea (es. \doublespacing)
\usepackage{geometry}             % Serve per personalizzare i margini e le dimensioni della pagina
\usepackage{enumitem}             % Dà il controllo totale su elenchi puntati e numerati
\usepackage{hyperref}             % Trasforma i riferimenti incrociati, URL e l'indice in link cliccabili nel PDF

\geometry{margin=1in}

\title{Progetto BDraw}
\author{Mattia Zanini}
\date{}

\begin{document}

\maketitle
\doublespacing

\section{Introduzione}

Il celebre problema della brachistocrona, ovvero la ricerca della curva che minimizza il tempo di percorrenza
tra due punti sotto l'effetto della gravità, rappresenta una delle pietre miliari della storia della matematica.
La sua soluzione, un arco di cicloide, fu scoperta oltre tre secoli fa, ma il fascino di questa sfida risiede
ancora oggi non solo nel risultato, quanto nelle geniali intuizioni e negli strumenti matematici che la sua
ricerca ha saputo generare.

L'idea del progetto nasce dalla volontà di far sperimentare in prima persona ciò che per i matematici del
Seicento rappresentava una vera competizione intellettuale tra menti brillanti. L'obiettivo è duplice: da un
lato rendere tangibile l'evoluzione degli strumenti computazionali, mostrando come sfide che impegnarono i geni
di tutta Europa siano oggi accessibili a chiunque grazie ai computer; dall'altro offrire un'esperienza divertente
dove si può cercare la curva ottimale per via intuitiva, ``andando a occhio''.

Il software trasforma così un problema puramente teorico in un'attività pratica e interattiva: l'utente può
disegnare curve arbitrarie e confrontare direttamente il tempo ottenuto con quello della soluzione ottimale,
apprezzando visivamente e quantitativamente le differenze tra un percorso generico e la curva ideale.

Per realizzare questo obiettivo, l'applicazione mette a disposizione un'area di disegno dove tracciare il
percorso e osservare il moto di una pallina lungo la rotta scelta. La velocità di percorrenza è mostrata sia
attraverso una simulazione fisica in tempo reale, sia mediante il calcolo teorico preciso del moto sulla
spezzata disegnata. Le curve possono essere fornite sia graficamente (tramite trascinamento del mouse) sia
numericamente (tramite equazioni matematiche) \textcolor{darkred}{(DA IMPLEMENTARE FORSE ALLA FINE)},
alimentando lo stesso motore di calcolo sottostante.

\subsection{Storia della Brachistocrona}

Il problema della brachistocrona, dalla parola greca composta \textit{brákhistos} (\foreignlanguage{greek}
{βράχιστος}, ``il più breve'') e \textit{khrónos} (\foreignlanguage{greek}{χρόνος}, ``tempo''), rappresenta
una delle pietre miliari della storia della matematica e del calcolo delle variazioni.

\subsection{La formulazione di Galileo (1638)}

Il problema fu formulato per la prima volta da Galileo Galilei nel 1638. Nel suo studio sul moto dei corpi,
Galileo si chiese quale fosse la curva lungo la quale un punto materiale, scorrendo senza attrito sotto
l'azione della gravità costante, impiegasse il tempo minimo per percorrere la distanza tra due punti fissati.
Per motivi estetici e per la mancanza di strumenti matematici appropriati, dato che il calcolo infinitesimale
era ancora agli albori, Galileo congetturò erroneamente che la soluzione fosse un arco di circonferenza.

Questa ipotesi, sebbene sbagliata, non era irragionevole: come mostrato da Buttazzo e Mintchev, l'errore di
valutazione di Galileo era sorprendentemente contenuto. In particolare, considerando il percorso necessario
per tornare alla stessa altezza iniziale ($H=0$), il tempo di percorrenza lungo la cicloide è circa il
\textbf{96\%} di quello impiegato lungo una semicirconferenza. Questo rende l'intuizione di Galileo molto
vicina alla soluzione ottimale dal punto di vista numerico, nonostante l'errore concettuale sulla forma
della curva.

\subsection{La sfida di Johann Bernoulli (1696)}

La soluzione corretta arrivò quasi sessant'anni dopo, quando nel giugno 1696 Johann Bernoulli lanciò una sfida
pubblica ai matematici di tutta Europa, pubblicando il problema sulla rivista \textit{Acta Eruditorum} di Lipsia.
La risposta non si fece attendere: entro la fine del 1697, soluzioni furono proposte da alcuni dei più grandi
studiosi dell'epoca:

\begin{itemize}
    \item \textbf{Isaac Newton}, che secondo la tradizione avrebbe risolto il problema in una sola notte,
          dopo averlo ricevuto per posta
    \item \textbf{Gottfried Wilhelm Leibniz}, co-fondatore del calcolo infinitesimale
    \item \textbf{Jacob Bernoulli}, fratello di Johann, che introdusse un metodo più generale
    \item Lo stesso \textbf{Johann Bernoulli}, con un approccio basato sul principio di Fermat
    \item \textbf{Guillaume de l'Hôpital}, autore del celebre trattato sulle coniche
\end{itemize}

Ognuno di questi matematici utilizzò metodi diversi, ma il vero valore di questa competizione intellettuale
non fu solo la risposta in sé. La brachistocrona è storicamente significativa perché \textbf{per risolvere
    questo problema furono sviluppati strumenti matematici che si sarebbero rivelati fondamentali in moltissimi
    altri contesti}. La necessità di trovare una curva che minimizzasse un funzionale (e non semplicemente un
numero) richiese un salto concettuale che portò alla nascita del calcolo delle variazioni.

\subsection{La soluzione: un arco di cicloide}

La soluzione si rivelò essere un \textbf{arco di cicloide}, la curva descritta da un punto fisso su una
circonferenza che rotola senza strisciare lungo una retta. La soluzione in forma parametrica è:

\[
    \begin{cases}
        x(t) = c(t - \sin t) \\
        y(t) = c(1 - \cos t)
    \end{cases}
    \qquad t \in [0, \tau]
\]

dove la costante $c$ e il parametro $\tau$ sono determinati dalle condizioni al contorno.

\subsection{Nascita del calcolo delle variazioni}

Il problema della brachistocrona è probabilmente il primo problema di calcolo delle variazioni in
dimensione infinita. La formulazione matematica richiede di minimizzare:

\[
    T(u) = \frac{1}{\sqrt{2g}} \int_{0}^{L} \sqrt{\frac{1 + |u'(x)|^2}{u(x)}} \, dx
\]

che rappresenta il tempo totale di percorrenza. Dall'equazione di Eulero-Lagrange (nella forma integrata
di DuBois-Reymond) si ottiene:

\[
    (1 + |u'(x)|^2) \ u(x) = 2c
\]

Questa equazione differenziale non lineare portò alla soluzione cicloidale. Il tempo di percorrenza lungo
la cicloide risulta circa il 16\% inferiore rispetto al percorso rettilineo tra gli stessi due punti.

L'importanza storica di questo risultato va ben oltre la risposta specifica: le tecniche sviluppate per
la brachistocrona gettarono le basi per intere nuove branche della matematica applicata, con applicazioni
che spaziano dalla fisica teorica all'ingegneria, dalla robotica alla teoria del controllo ottimo.

\section{Requisiti Funzionali}

\subsection{Input e Disegno}

\begin{itemize}
    \item L'utente deve poter disegnare una curva a mano libera sull'area di disegno tramite il trascinamento
          del mouse
    \item Il sistema deve permettere l'inserimento di curve tramite equazioni matematiche, oltre all'input
          grafico \textcolor{darkred}{(DA IMPLEMENTARE FORSE ALLA FINE)}
    \item Il sistema deve poter campionare correttamente la posizione del mouse sull'area di disegno
    \item Il sistema deve garantire la monotonia crescente della curva rispetto all'asse X, rimuovendo i
          punti che violano questo vincolo
    \item Il sistema deve unire tutte le posizioni campionate per formare una curva continua
    \item Il sistema deve permettere di pulire l'area di disegno per creare una nuova curva
\end{itemize}

\subsection{Curve Predefinite e Calcolo Teorico}

\begin{itemize}
    \item Il sistema deve poter generare percorsi predefiniti: linea retta e cicloide
    \item Il sistema deve calcolare il tempo teorico di percorrenza sulla spezzata disegnata
\end{itemize}

\subsection{Simulazione Fisica}

\begin{itemize}
    \item Il sistema deve eseguire una simulazione fisica della pallina che percorre la curva sotto
          l'effetto della gravità
    \item La velocità della pallina deve variare lungo il percorso in base alla quota, secondo il principio
          di conservazione dell'energia
    \item Il sistema deve misurare e visualizzare il tempo effettivo di percorrenza durante la simulazione
    \item Il sistema deve permettere di ripetere la simulazione sulla stessa curva
\end{itemize}

\subsection{Visualizzazione e Confronto}

\begin{itemize}
    \item Il sistema deve visualizzare il confronto tra tempo teorico e tempo effettivo di percorrenza
    \item Il sistema deve loggare informazioni di debug su console (parametri fisici, stato della simulazione)
\end{itemize}

\section{Requisiti Non Funzionali}

\subsection{Prestazioni, Responsività e Logging}

\begin{itemize}
    \item Il loop di simulazione deve essere eseguito con una frequenza tale da garantire fluidità, per
          mantenere la finestra dell'applicazione responsiva.
    \item Il sistema deve fornire informazioni di debug su console tramite la libreria \texttt{spdlog}
          configurata adeguatamente
\end{itemize}

\subsection{Calcolo e Simulazione Fisica}

\begin{itemize}
    \item Il tempo teorico è calcolato in modo esatto sulla spezzata: poiché la curva è costante a tratti
          su ogni segmento, la formula di Galileo fornisce il valore preciso senza approssimazioni numeriche
    \item La posizione della pallina in ciascun frame deve essere determinata collegando il tempo fisico
          reale con il fluire dei frame, garantendo coerenza fisica nella visualizzazione
\end{itemize}

\subsection{Elaborazione e Rendering}

\begin{itemize}
    \item Per garantire la monotonia crescente della curva viene eseguito un postprocessing sui punti
          campionati al momento del rilascio del mouse (terminazione del disegno)
    \item La curva disegnata deve apparire continua all'utente, unendo tutte le posizioni campionate dal mouse
    \item La pallina deve rimanere visivamente sopra la curva senza compenetrazioni grafiche, calcolando
          la normale della superficie nel punto di contatto
\end{itemize}

\section{Requisiti di Sistema}

\begin{itemize}
    \item \textbf{Linguaggio}: C++17
    \item \textbf{Framework UI}: Qt 6.10
    \item \textbf{Calcolo Numerico}: Armadillo
    \item \textbf{Logging}: \texttt{spdlog}
    \item \textbf{Build System}: CMake 3.16+
\end{itemize}

\end{document}
